\documentclass[11pt]{article}
\usepackage[margin=0.85in]{geometry}
\usepackage{booktabs}
\usepackage{float}
\usepackage{graphicx}
\usepackage{caption}
\usepackage{array}
\usepackage{amsmath, amssymb}
\usepackage{hyperref}
\usepackage{xcolor}
\usepackage{listings}
\lstset{
  basicstyle=\ttfamily\footnotesize,
  breaklines=true,
  frame=single,
  framesep=4pt,
  xleftmargin=4pt,
  xrightmargin=4pt,
  backgroundcolor=\color{black!3},
}

\title{Two PSID Objects on the Down-Payment Constraint}
\author{}
\date{Built May 26, 2026}

\begin{document}
\maketitle

\noindent This note documents two PSID empirical objects that together discipline the collateral constraint in the model and provide reduced-form evidence that the constraint matters for fertility. The model object is
\[
b_t \;\geq\; (1-\phi)\, p_i\, H_k,
\]
with calibrated $\phi = 0.80$, so 20\% of the price of an owned home in liquid wealth. Ownership in the model unlocks family-size housing services $\bar h(1) > \bar h(0)$. The model therefore predicts (i) actual first-time buyers cluster at or below the 20\% threshold, and (ii) renters whose liquid wealth falls short of the implied down payment for a typical owned home are disproportionately childless. The two objects below test these two predictions.

\section{Object A: Bunching at the Collateral Frontier}

\subsection*{Construction}

For each PSID head observed in $1984$--$2019$, identify the first survey year in which the head is an owner ($\texttt{HOMEOWN} = 1$) and was previously observed as a renter ($\texttt{HOMEOWN} = 2$). Require $\texttt{YEARMOVED\_} \in [t-2, t]$ so reported value and mortgage are at or near origination. The down-payment share at first purchase is
\[
\frac{\mathrm{DP}}{P} \;=\; 1 \;-\; \frac{\texttt{HOMEMORTOTR}_t}{\texttt{HOMEVALUER}_t}.
\]
Cash buyers ($\texttt{HOMEMORTOTR} = 0$) are 17.9\% of the sample and are shown separately. The headline density is plotted conditional on having a mortgage so the cash spike at $\mathrm{DP}/P = 1$ does not visually distract from the bunching story.

\subsection*{Result}

\begin{figure}[H]
\centering
\includegraphics[width=0.85\textwidth]{\detokenize{../code/data/psid_followup_mar2026/output/dp_at_first_purchase_v1/dp_share_density_nocash_v1.png}}
\caption{Kernel density of the down-payment share at first home purchase, PSID first-time buyers with a mortgage, $1984$--$2019$. Dotted lines mark $3.5\%$ (FHA), $5\%$, and $10\%$ conventional floors; the dashed red line is the $20\%$ GSE conforming threshold (the calibrated $1-\phi$).}
\label{fig:dp_density}
\end{figure}

\begin{table}[H]
\centering
\caption{Down-payment share at first purchase, mortgaged buyers.}
\begin{tabular}{lr}
\toprule
Statistic & Value \\
\midrule
First-time mortgaged buyers ($N$) & $252$ \\
Median $\mathrm{DP}/P$ & $0.16$ \\
Share with $\mathrm{DP}/P \leq 0.035$ (FHA) & $0.241$ \\
Share with $\mathrm{DP}/P \leq 0.05$            & $0.282$ \\
Share with $\mathrm{DP}/P \leq 0.10$            & $0.378$ \\
Share with $\mathrm{DP}/P \leq 0.20$ (GSE)      & $\mathbf{0.566}$ \\
Cash buyers (separate)                          & $0.179$ \\
\bottomrule
\end{tabular}
\end{table}

\noindent The density is monotone decreasing from $0$ to $0.20$ with a visible step right at the $20\%$ GSE line and a thin right tail. The bunching is not a $2008$--$2010$ crisis artifact: the share at or below $20\%$ is $47\%$ in $1995$--$2004$, $43\%$ in $2005$--$2010$, and $50\%$ in $2011$--$2019$ (raw shares including cash buyers; the mortgaged-only versions are higher).

\subsection*{What Object A buys}

Direct, model-free evidence that real first-time buyers operate at the collateral frontier the model writes down. The DP constraint is where the modal mortgaged first-time buyer lands.

\section{Object B: DP Shortfall Predicts Childlessness}

\subsection*{Construction}

Build a per-year target home price as the median value among observed first-time PSID buyers in the same year-bin (Table~\ref{tab:p_target}). For each PSID head observed as a renter at ages $25$--$30$, define
\[
\text{shortfall}_{it} \;=\; \mathbf{1}\!\left\{ b_{it} \;<\; 0.20 \cdot P^{\mathrm{target}}_{t} \right\},
\]
where $b_{it} = \texttt{NETWORTH2R}$ (PSID net worth excluding home equity). Collapse to one observation per ID (mean over the $25$--$30$ window), keep observations followed to age $45$, and compare childlessness by $45$ across shortfall status within income tercile.

\begin{table}[H]
\centering
\caption{Target home price $P^{\mathrm{target}}_t$ by year-bin (PSID first-time buyer median home value, real).}
\label{tab:p_target}
\begin{tabular}{lrr}
\toprule
Year bin & $P^{\mathrm{target}}$ (real \$) & First-time buyers \\
\midrule
$1995$--$2004$ & $179{,}300$ & $95$ \\
$2005$--$2010$ & $188{,}844$ & $128$ \\
$2011$--$2019$ & $218{,}775$ & $169$ \\
\bottomrule
\end{tabular}
\end{table}

\subsection*{Result}

\begin{figure}[H]
\centering
\includegraphics[width=0.78\textwidth]{\detokenize{../code/data/psid_followup_mar2026/output/dp_shortfall_v1/childless_by_shortfall_inc_tercile_v1.png}}
\caption{Share childless by 45 among PSID renters at ages $25$--$30$, split by DP-shortfall status within income tercile.}
\label{fig:shortfall_bars}
\end{figure}

\begin{table}[H]
\centering
\caption{Childlessness by $45$ by shortfall status and income tercile (PSID renters at $25$--$30$).}
\begin{tabular}{lrrrrr}
\toprule
Income tercile & DP sufficient & $N$ & DP shortfall & $N$ & Gap \\
\midrule
Bottom & $0.227$ & $19$    & $0.250$ & $1{,}237$ & $+2.3$ pp \\
Middle & $0.145$ & $37$    & $0.274$ & $864$     & $+\mathbf{12.9}$ pp \\
Top    & $0.103$ & $59$    & $0.230$ & $758$     & $+\mathbf{12.7}$ pp \\
\bottomrule
\end{tabular}
\end{table}

\begin{table}[H]
\centering
\caption{OLS of childless-by-$45$ on DP shortfall and controls. PSID renters at $25$--$30$, pre-first-birth, $N = 2{,}961$. PSID weights, robust SEs.}
\begin{tabular}{lrrr}
\toprule
Specification & $\hat\beta_{\text{shortfall}}$ & SE & $p$ \\
\midrule
Shortfall ($20\%$)                          & $0.115$ & $0.039$ & $0.003$ \\
Shortfall ($20\%$) $+$ $\log$ income        & $0.116$ & $0.040$ & $0.003$ \\
Shortfall ($20\%$) $+$ $\log$ income $+$ $\log$ NW & $0.122$ & $0.041$ & $0.003$ \\
Strict shortfall ($10\%$) $+$ $\log$ income & $0.068$ & $0.036$ & $0.063$ \\
\bottomrule
\end{tabular}
\end{table}

\noindent DP-shortfall renters are about $12$ pp more childless by $45$, controlling for log income and log net worth. The $10\%$ strict threshold is weaker, consistent with $20\%$ being the binding margin in both this object and the calibrated $\phi$.

\subsection*{What Object B buys}

Direct reduced-form evidence that the constraint matters for fertility, not just for tenure. The threshold is sharper than raw wealth because it varies with both wealth and prices (the threshold $0.20 \cdot P^{\mathrm{target}}_t$ moves over time), and the within-income split nets out the income channel.

\section{Why Both Together Make the Case}

Object~A alone says the constraint binds for buyers; a critic can respond that non-buyers stay renters for reasons unrelated to fertility. Object~B alone says renters with little wealth have fewer kids; a critic can respond that the result is just income or wealth. Together, the \emph{same} $20\%$ threshold that buyers bunch at also separates childless from parent renters within income tercile. Two complementary moments of the same distribution, both consistent with the calibrated $\phi$.

\section{Draft Pitch Frames}

These are drafts to paste into \texttt{latex/may\_29\_project\_presentation.tex}. Style matches the existing PSID frames around lines $1188$ and $1235$. Adjust hyperref targets if needed.

\subsection*{Frame 1: Bunching at the collateral frontier}

\begin{lstlisting}[language={}]
\begin{frame}{Buyers Bunch at the 20\% Down-Payment Constraint}
\begin{columns}[T]
\begin{column}{0.58\textwidth}
\centering
\includegraphics[width=\linewidth]{../code/data/psid_followup_mar2026/output/dp_at_first_purchase_v1/dp_share_density_nocash_v1.png}
\end{column}
\begin{column}{0.40\textwidth}
\small
PSID first-time buyers with a mortgage, $1984$--$2019$ ($N=252$).
\medskip
\begin{itemize}
  \item $\mathrm{DP}/P = 1 - \text{mortgage}/\text{value}$ at first owner-year.
  \item Density is monotone from $0$ to $0.20$ with a visible step at the GSE $20\%$ line.
  \item $\mathbf{57\%}$ of mortgaged buyers put down $\leq 20\%$; $38\%$ put down $\leq 10\%$.
  \item Stable across $1995$--$2004$, $2005$--$2010$, $2011$--$2019$: not a crisis artifact.
\end{itemize}
\end{column}
\end{columns}

\medskip
\footnotesize The calibrated $\phi = 0.80$ matches the threshold the modal mortgaged first-time buyer hits in the data.
\end{frame}
\end{lstlisting}

\subsection*{Frame 2: Down-payment shortfall predicts childlessness}

\begin{lstlisting}[language={}]
\begin{frame}{Renters Below the Down-Payment Frontier Stay Childless}
\begin{columns}[T]
\begin{column}{0.58\textwidth}
\centering
\includegraphics[width=\linewidth]{../code/data/psid_followup_mar2026/output/dp_shortfall_v1/childless_by_shortfall_inc_tercile_v1.png}
\end{column}
\begin{column}{0.40\textwidth}
\small
PSID renters at $25$--$30$, pre-first-birth, $N=2{,}961$.
\medskip
\begin{itemize}
  \item Shortfall $=$ liquid wealth $< 0.20 \times$ median first-time-buyer home value in that year-bin.
  \item Within \emph{every} income tercile, shortfall renters are more childless by $45$.
  \item Mid- and top-income terciles: gap is $\mathbf{\approx 13}$ pp.
  \item Regression coefficient $0.122$ (SE $0.041$, $p=0.003$), controlling for $\log$ income and $\log$ net worth.
\end{itemize}
\end{column}
\end{columns}

\medskip
\footnotesize The 20\% threshold is the same number that buyers bunch at and the same number the model calibrates as $1 - \phi$.
\end{frame}
\end{lstlisting}

\section{Caveats}

\begin{itemize}
  \item \textbf{Sample sizes}: PSID is small. Object~A has $N = 252$ mortgaged first-time buyers, Object~B has $N = 2{,}961$ pre-birth renters followed to $45$. Within bottom income tercile, the DP-sufficient comparison group in B has only $N = 19$, so that row should be read as descriptive only.
  \item \textbf{Origination proxy}: $\mathrm{DP}/P$ is computed from current $\texttt{HOMEVALUER}$ and current $\texttt{HOMEMORTOTR}$ at the first surveyed owner-year, with $\texttt{YEARMOVED\_} \in [t-2, t]$. Some drift from true origination is unavoidable. House-price appreciation post-purchase biases reported $\mathrm{DP}/P$ \emph{up}, so the bunching headline is conservative.
  \item \textbf{Gifts and transfers}: A buyer financed by a family gift looks like a high-DP buyer; this also biases the bunching estimate down. The true share at the collateral floor is plausibly higher than reported.
  \item \textbf{Target price in B}: uses national year-bin median first-time-buyer value. A state- or metro-level analog would tighten the threshold; the PSID shelf has $\texttt{STATEFIPS\_}$ but not MSA. State-level version is straightforward as a robustness pass.
  \item \textbf{Outright cash buyers} are $17.9\%$ of A's sample. They are excluded from the headline density and reported separately; the all-buyers density and summary are still saved as $\texttt{dp\_share\_density\_v1.png}$ and $\texttt{dp\_share\_bunching\_summary\_v1.csv}$.
\end{itemize}

\section{Code and Artifacts}

\begin{itemize}
  \item Do-files:
  \texttt{code/data/psid\_followup\_mar2026/dp\_at\_first\_purchase\_v1.do} (Object~A),
  \texttt{code/data/psid\_followup\_mar2026/dp\_shortfall\_v1.do} (Object~B).
  \item Outputs:
  \texttt{code/data/psid\_followup\_mar2026/output/dp\_at\_first\_purchase\_v1/} and
  \texttt{code/data/psid\_followup\_mar2026/output/dp\_shortfall\_v1/}.
\end{itemize}

\end{document}
